\chapter{Introduction}
\label{ch:intro}
% Chapter 1: Introduction (8-10 pages)

%
% Section: Motivation
%
\section{Motivation}
\label{sec:intro:motivation}
\graffito{Leishmaniasis is a neglected tropical disease affecting millions worldwide.}
% Content: Leishmaniasis as NTD, diagnostic challenges in resource-limited settings,
% gap between flagship LLMs and clinical deployment reality,
% need for knowledge augmentation in small language models.

%
% Section: Problem Statement
%
\section{Problem Statement}
\label{sec:intro:problem}
% Content: Limitations of zero-shot LLMs for rare disease diagnosis,
% data scarcity in specialized medical domains,
% multimodal nature of clinical diagnosis (text + images).

%
% Section: Research Questions
%
\section{Research Questions}
\label{sec:intro:rq}
% RQ1: Can RAG improve diagnosis accuracy of small language models for Leishmaniasis?
% RQ2: How does multimodal retrieval (text + clinical images) compare to text-only retrieval?
% RQ3: Which retrieval architecture (dense, sparse, hybrid) performs best for medical case retrieval?

%
% Section: Contributions
%
\section{Contributions}
\label{sec:intro:contributions}
% - Two-lane multimodal RAG architecture (Lane 1: Text, Lane 2: Vision-Language)
% - Curated Leishmaniasis multimodal dataset (143 verified cases)
% - Comprehensive evaluation framework with RAGAS metrics
% - Analysis of RAG vs No-RAG across model sizes

%
% Section: Thesis Structure
%
\section{Thesis Structure}
\label{sec:intro:structure}
% Overview of remaining chapters
